\documentclass[11pt]{article}
\usepackage{manual_docs_style}
\externaldocument{build/dataset_manifest_schema}[dataset_manifest_schema.pdf]
\externaldocument{build/candidate_generation}[candidate_generation.pdf]
\externaldocument{build/cheap_filter_metrics_stage}[cheap_filter_metrics_stage.pdf]
\externaldocument{build/surrogate_confidence_scoring_stage}[surrogate_confidence_scoring_stage.pdf]

\begin{document}

\docTitle{Stage: Generate Question}
\phantomsection\label{docstart:generate_question_stage}

This stage exports selected simulator runs into the canonical TraceBench question bundle: the standard exchange format consumed by agentic rollout, baseline evaluation, and downstream analysis.

\section*{Required Inputs}

For each \termSimulator, question generation requires:

\begin{itemize}
\item \code{datasets/<dataset_id>/sample_manifest.json}: selected train/test panels.
\item \code{datasets/<dataset_id>/selected_runs.jsonl}: selected runs and split metadata.
\item \code{datasets/<dataset_id>/selection_report.json}: frozen selection diagnostics.
\item \code{runs/<run_id>/data.parquet}: materialized time-series artifacts.
\item \code{runs/model_record.json}: runtime inventory for selected runs.
\item \code{plans/<production_policy_id>/run_nodes.jsonl}: canonical production run recipes and hashes.
\item \code{plans/<production_policy_id>/run_edges.jsonl}: production baseline/intervention/time-zero relationships.
\item \code{experiment_config.json}
\item \code{description_levels.json}
\item \code{noise_adder.py}
\end{itemize}

Question generation itself should not need to rescore the whole candidate pool.
The preferred design writes the selected dataset artifacts first.

\section*{Step 1: Dataset Selection}

The selection step builds a selected dataset manifest.
It is documented in the Dataset Manifest Schema and is conceptually upstream of this stage.

Recommended invocation:

\begin{DocCode}
env/bin/python workflows/generate/select_dataset.py \
  --model SIMULATOR \
  --policy PRODUCTION_POLICY_ID \
  --dataset DATASET_ID \
  --surrogate-id SURROGATE_ID \
  [--row-slug ROW_SLUG [ROW_SLUG ...]]
\end{DocCode}

\section*{Step 2: Generate Exported Questions}

This step is deterministic once \code{sample_manifest.json} exists.

\begin{DocCode}
env/bin/python workflows/generate/create_exam_questions_tracebench_cls_from_registry.py \
  --model SIMULATOR [SIMULATOR ...] \
  --dataset DATASET_ID \
  --row-slug SLUG [SLUG ...]
\end{DocCode}

It materializes the canonical exported bundle as:
\begin{DocCode}
  TraceBench_questions/<simulator_id>/
    questions.json
    dataframes/
      `-- *.parquet
    noise_adder.py
    sample_manifest.json
    model_record.json
\end{DocCode}
Where:
\section*{\texorpdfstring{\titlecode{sample_manifest.json}}{sample\_manifest.json}}

The manifest is keyed by \code{shot_slug}.
Each \code{shot_slug} maps to a non-empty list of entries, one per resolved \code{(shot_slug, test_set_slug, seed)} combination.

Each entry contains:

\begin{itemize}
\item \code{train_samples_baselines}
\item \code{train_samples}
\item \code{train_samples_hashes}
\item \code{test_samples_baselines}
\item \code{test_samples}
\item \code{test_samples_hashes}
\item \code{other_samples}
\item \code{seed}
\item \code{test_set_slug}
\item \code{shot_slug_recipe}
\item \code{accuracy_with_baselines}, when adversarial selection is evaluated
\item \code{accuracy_with_baselines_all_samples}, when evaluated
\item \code{train_test_sample_hash}
\end{itemize}

The pair \code{(test_set_slug, seed)} has a one-to-one mapping to \code{test_samples_hashes}.

\section*{\texorpdfstring{\titlecode{questions.json}}{questions.json}}

The exported \code{questions.json} contains:

\begin{itemize}
\item \code{version}
\item \code{questions}
\item \code{label_int_mapping}
\item \code{ground_truth_information}
\item \code{environment_name}
\end{itemize}

\code{ground_truth_information} is derived from the selected production run graph and selected dataset manifest.
For each selected sample, it records initial parameters, changed parameter, intervention time, new value, and \code{first_diff} values.
No-intervention samples use the no-intervention class and null intervention fields.
Surrogate confidence values may be copied into hidden provenance metadata, but they must not appear in the agent-visible prompt or files.

\section*{Question Text}
\phantomsection\label{sec:generate-question-question-text}

\code{question_text} contains the structured components used to render the final agent prompt.
The final prompt is produced by \code{shared/prompts.py::render_tracebench_agent_prompt(question_text:Mapping[str, Any]) -> str}.

The payload contains:
\begin{itemize}
\item \code{sample_source}: description of which sample folders are available.
\item \code{environment_description}: environment description controlled by \termDescLevel.
\item \code{observed_columns}: observed-column and signal descriptions.
\item \code{intervention_semantics}: shared intervention semantics.
\item \code{allowed_labels}: canonical allowed labels.
\item \code{label_space}: rendered text form of the allowed labels.
\item \code{no_change_guidance}: guidance for the no-change label.
\item \code{fewshot_context}: few-shot context when train samples are present.
\item \code{task_artifact}: task artifact instructions.
\item \code{prediction_format}: required prediction/output format.
\item \code{mode_specific_requirements}: requirements specific to direct or code mode.
\item \code{evaluation}: evaluation text.
\item \code{runtime_constraints}: runtime constraints.
\item \code{ordered_field_agent_prompt}: canonical rendering order and per-field separators. It looks like: 
\begin{DocCode}
    "ordered_field_agent_prompt": [
    ["sample_source", " "],
    ["environment_description", "\n\n"],
    ["observed_columns", "\n\n"],
    ["intervention_semantics", "\n\n"],
    ["label_space", "\n\n"],
    ["no_change_guidance", "\n\n"],
    ["task_artifact", "\n"],
    ["prediction_format", "\n\n"],
    ["fewshot_context", "\n\n"],
    ["mode_specific_requirements", "\n\n"],
    ["evaluation", "\n\n"],
    ["runtime_constraints", ""]
  ]
\end{DocCode}
\end{itemize}

The renderer renders the fields listed in \code{question_text.ordered_field_agent_prompt} in order.
Each entry specify the separator to use before the next rendered field.
Missing or empty fields are skipped, and placeholder references inside field values are resolved before concatenation.

These values are resolved from \code{shared/config/tracebench_shared_description.json}, \code{description_levels.json}, and the selected dataset manifest.
The function \code{render_tracebench_agent_prompt(...)} renders the final agent question from these entries using \code{ordered_field_agent_prompt}.


\section*{Validation Rules}

\begin{itemize}
\item Question generation must not select new samples.
\item It may only export samples already present in \code{datasets/<dataset_id>/selected_runs.jsonl} and grouped by \code{sample_manifest.json}.
\item Every selected \code{run_id} must exist in \code{run_nodes.jsonl} and \code{model_record.json}.
\item Every selected run must have a materialized \code{data.parquet} artifact.
\item Test-set length must match the requested \code{number_test_samples}.
\item Test-set class counts should be balanced, with any two class counts differing by at most one.
\item Train-set length must equal \code{number_train_samples_per_class * num_labels}.
\item Train and test samples must not overlap within a question.
\item Comparable rows sharing \code{test_set_slug} and seed must share the same underlying test set.
\item Exported question hashes must be derived from selected sample hashes.
\end{itemize}

\end{document}
